% ======================================================================
%  sec00.tex —— 第 0 节 引言：为什么需要概形？
%  底本：原书 pp.5–8（印刷页 = PDF 页）
%  体例：每节开头 \section + \subsection 与原文编号一致；
%        定理类环境紧跟 \label{类型:编号}；译注用 yizhu 环境。
% ======================================================================

\section{引言：为什么需要概形？}\label{sec:0}

% ---------------------------------------------------------------- p.5
\subsection{代数集}\label{sec:0.1}

在概形理论之前，代数几何关注的对象是代数集。

\begin{definition}\label{def:0.1.1}
设 $k$ 为代数闭域。
\begin{enumerate}[label=(\roman*),leftmargin=2em,itemsep=0.25ex]
  \item 仿射空间 $\A^n_k$ 上的 \emph{Zariski 拓扑}，是指以如下子集为基的拓扑：这些子集等于某个多项式 $f \in k[x_1,\dots,x_n]$ 的非零点轨迹 $U(f)$。
  \item $\A^n_k$ 中的一个\emph{嵌入仿射代数集}\footnote{有人使用 variety（簇）一词，也有人把它留给不可约代数集。}，是指 Zariski 拓扑下的闭子空间。
  \item 一个\emph{嵌入拟仿射代数集}，是指某个嵌入仿射代数集的 Zariski 开子集。
\end{enumerate}
\end{definition}

\begin{yizhu}
「嵌入的」（embedded）三个字是原书的用意所在：这里的一切对象都先验地放在某个固定的 $\A^n_k$ 中，闭集由环境空间的多项式方程定义。nLab 的「Zariski topology」条目把这一构造说成一部词典：闭子集对应多项式环中的理想（取 $V(I)$），不可约闭子集对应素理想，而闭点对应极大理想。词典的最后一条正是第 \ref{sec:0.2} 小节从 $\Spm$ 走向 $\Spec$ 的伏笔：一旦把「闭点」换成「素理想」，环境空间就不再必要。
[来源：nLab "Zariski topology"，\url{https://ncatlab.org/nlab/show/Zariski+topology}（Last revised 2024-12-19；2026-11-17 访问）。]
\end{yizhu}

\begin{example}\label{ex:0.1.2}
$\A^n_k$ 的任意有限子集都是嵌入仿射代数集。
\end{example}

\begin{example}\label{ex:0.1.3}
$\Z$ 不是 $\A^1_\C$ 中的嵌入仿射代数集。
\end{example}

类似地，用齐次多项式与射影空间 $\PP^n_k$，可以定义嵌入（拟）射影代数集。

上述定义有大量缺陷。一个明显的缺陷是：「嵌入的」代数集这一概念并非内蕴的。

\begin{example}\label{ex:0.1.4}
嵌入仿射代数集 $\A^1_k \subseteq \A^1_k$ 与 $\A^1_k \subseteq \A^2_k$ 应当被视为同一个代数集。
\end{example}

\begin{notation}\label{not:0.1.5}
为弥补这一点，我们需要若干记号。
\begin{enumerate}[label=(\roman*),leftmargin=2em,itemsep=0.25ex]
  \item 对理想 $I \subseteq k[x_1,\dots,x_n]$，记 $Z(I) \subseteq \A^n_k$ 为 $I$ 中多项式的公共零点之轨迹。
  \item 对 Zariski 闭子集 $X \subseteq \A^n_k$，记 $I(X) \subseteq k[x_1,\dots,x_n]$ 为在 $X$ 上消失的全部多项式构成的理想。
\end{enumerate}
\end{notation}

回忆：理想 $I$ 称为\emph{根理想}，若 $I = \sqrt{I}$。

\begin{theorem}[Hilbert 零点定理]\label{thm:0.1.6}
存在双射
\[
\begin{array}{rcl}
\{\,k[x_1,\dots,x_n] \text{ 的根理想}\,\} & \longleftrightarrow & \{\,\A^n_k \text{ 的 Zariski 闭子集}\,\}\\[0.5ex]
I & \longmapsto & Z(I)\\
I(X) & \longmapsfrom & X.
\end{array}
\]
\end{theorem}

定理的一部分断言：$\A^n_k$ 的点集与 $k[x_1,\dots,x_n]$ 的极大理想集一一对应。作为推论，$Z(I)$ 与包含 $I$ 的极大理想一一对应；后者又可进一步等同于 $R := k[x_1,\dots,x_n]/I$ 的极大理想。

注意 $I$ 是根理想当且仅当 $R$ 是既约的，即不含幂零元。这为下述定义提供了依据。

\begin{definition}\label{def:0.1.7}
\emph{仿射代数 $k$-集}，是指某个有限生成（交换含幺）既约 $k$-代数 $R$ 的\emph{极大谱} $\Spm R$（即极大理想之集）。我们在其上赋予 Zariski 拓扑，其开集基由
\[
U(f) := \{\,\mathfrak m \in \Spm R \mid f \notin \mathfrak m\,\}, \qquad f \in R,
\]
给出。
\end{definition}

\begin{yizhu}
零点定理的完整形式比原书陈述的双射更强：它同时断言，对任意极大理想 $\mathfrak m \subset k[x_1,\dots,x_n]$，剩余域扩张 $\kappa(\mathfrak m)/k$ 是有限的；并且任一根理想都等于包含它的极大理想之交。「交换环 $\to$ 仿射概形」这一侧之所以能够成立，靠的正是后一条：根理想被它在极大理想中的成员完全决定。Stacks Project 把这层内容单列为定理，并明确说明它可推广到任意有限生成 $k$-代数。
[来源：Stacks Project, Tag 00FV（Theorem 10.34.1，Hilbert Nullstellensatz），\url{https://stacks.math.columbia.edu/tag/00FV}（2026-11-17 访问）。]
\end{yizhu}

% ---------------------------------------------------------------- p.6
\begin{example}\label{ex:0.1.8}
$\Spm k[x] \simeq \A^1_k$。
\end{example}

代数与几何之间有如下对偶。
\[
\begin{array}{r|c}
\text{代数} & \text{几何}\\[0.4ex] \hline
\rule{0pt}{1.2em}\text{有限生成既约 $k$-代数}\,R & \text{仿射代数 $k$-集}\;X\\
\text{极大理想}\;\mathfrak m \subseteq R & \text{点}\;x \in X\\
\text{元素}\;f \in R & \text{函数}\;\varphi \colon X \to \A^1_k\\
\text{根理想}\;I \subseteq R & \text{Zariski 闭子集}\;Z \subseteq X.
\end{array}
\]

这里元素 $f \in R$ 对应于函数
\[
\varphi \colon \Spm R \to k, \qquad \mathfrak m \mapsto \overline{f},
\]
它把极大理想 $\mathfrak m$ 送到 $f$ 在 $\mathfrak m$ 的剩余域中的像 $\overline{f}$；该剩余域经由合成 $k \to R \to R/\mathfrak m$ 与 $\A^1_k$ 的底层集典范等同。

「对偶」一词意味着对应 $R \leftrightarrow X$ 是反变的。事实上，给定同态 $f \colon R' \to R$，我们得到连续映射
\[
\Spm R \to \Spm R', \qquad \mathfrak m \mapsto f^{-1}(\mathfrak m).
\]
但要注意：并非所有连续映射 $\Spm R \to \Spm R'$ 都由此得到，$R$ 也不被拓扑空间 $\Spm R$ 决定。

\begin{exercise}\label{exc:0.1.9}
证明 $\A^1_k \to \A^1_k$ 的任意双射对 Zariski 拓扑连续。找出其中来自环同态 $k[x] \to k[x]$ 的那些双射。
\end{exercise}

这引出下述定义。

\begin{definition}\label{def:0.1.10}
从 $\Spm R$ 到 $\Spm R'$ 的\emph{态射}，是指由同态 $R' \to R$ 给出的连续映射。
\end{definition}

于是，用态射把仿射代数 $k$-集粘起来，便可定义一般的代数 $k$-集；这正如人们用保持附加结构的映射，把带结构的欧氏空间粘成带结构的流形。

\begin{yizhu}
两张对照表都写作「Algbera」（p.6 起）、「Algera」（p.7 起），正文并有「Roughtly」「addiontal structrues」等拼写，均为原书笔误。译文按显然本意排印，不在表中保留误拼；读者若核校原文，请以原书为准。
[依据：原书 pp.6--8。]
\end{yizhu}

% ---------------------------------------------------------------- p.6--7
\subsection{缺陷}\label{sec:0.2}

代数 $k$-集的理论在代数与几何之间架起了一座桥。特别地，人们可以用拓扑与几何的方法（包括各种上同调理论）来研究有限生成既约 $k$-代数。然而，这些形容词对这座桥而言是不必要的限制。

第一，数论研究数域及其整数环，例如 $\Q$ 与 $\Z$。人们希望有与之对应的几何对象，因此愿意考虑一般的交换环而非 $k$-代数。但随即会意识到：极大谱 $\Spm$ 不够用。

\begin{example}\label{ex:0.2.1}
映射 $\Z \to \Q$ 不诱导映射 $\Spm \Q \to \Spm \Z$。因为 $(0) \subseteq \Q$ 在 $\Z$ 中的原像是非极大的素理想。
\end{example}

这提示我们：对一般的代数 $R$，应当考虑其\emph{素谱}，记作 $\Spec R$，而不只是极大谱。

第二，即使在研究有限生成代数时，也会自然遇到非有限生成的代数。

\begin{example}\label{ex:0.2.2}
设 $\mathfrak p \subseteq R$ 是有限生成代数的素理想。局部化 $R_{\mathfrak p}$ 及其完备化 $\widehat{R}_{\mathfrak p}$ 一般不是有限生成的。
\end{example}

\begin{yizhu}
$\Spm$ 的困难是结构性的：环同态 $R' \to R$ 只保证素理想的原像仍是素理想，却不保证极大理想的原像仍是极大理想，例 \ref{ex:0.2.1} 即 $\Z \to \Q$ 把 $(0)$ 拉回。换成 $\Spec$ 之后，「任意同态诱导连续映射」成为函子性的一部分，而不再是需要额外验证的条件；代价是把非闭点一并收入空间。原文第二张对照表末行「理想 $\leftrightarrow$ 闭子概形」正是这一代价换来的收益，见第 \ref{sec:0.3} 小节。
[依据：原书 pp.6--7；对照 Stacks Project, Tag 01IJ（Definition 26.9.1，概形与概形态射的定义），\url{https://stacks.math.columbia.edu/tag/01IJ}（2026-11-17 访问）。]
\end{yizhu}

当然，人们可以把注意力限制在诺特环上（从此过上幸福生活）。但我要反对这样一种错觉：人们关心的一切自然出现的环都是诺特的。

\begin{example}\label{ex:0.2.3}
诺特环对张量积不稳定：$\overline{\Q} \otimes_{\Q} \overline{\Q}$ 不是诺特环。
\end{example}

\begin{example}\label{ex:0.2.4}
$\Q$ 的阿代尔环不是诺特环。
\end{example}

\begin{example}\label{ex:0.2.5}
关于诺特对象的自然模问题，其本身未必是诺特的。我个人最喜欢的例子包括：形式联络的空间、形式群律的空间（栈）……
\end{example}

简言之，非诺特环有重要的应用，而本课程要处理的就是后者。

\begin{yizhu}
阿代尔环是 $\Q$ 的全部完备化（各 $\Q_p$ 与 $\R$）按「几乎处处取整数」的条件所作的限制积。这一构造把局部信息同时收进一个环，因而天然带有无穷多个因子；正因如此，它不满足升链条件。形式联络与形式群律的模空间同理：它们记录的是「在无穷远处逐阶给定」的数据，而非有限维的线性数据。
[来源：nLab "ring of adeles"，\url{https://ncatlab.org/nlab/show/ring+of+adeles}（Last revised 2025-07-01；2026-11-17 访问）。]
\end{yizhu}

最后，我们希望去掉既约性的限制。

\begin{example}\label{ex:0.2.6}
既约环对张量积不稳定：$k[x] \otimes_{k[x,y]} k[x,y]/(y-x^2)$ 不是既约环。几何上，这意味着 $Z(y)$ 与 $Z(y-x^2)$ 在 $\A^2_k$ 中不横截相交。
\end{example}

人们可能注意到：若去掉既约性，就相应地应考虑所有理想而非仅根理想；但此时构造 $I \mapsto Z(I)$ 不再是双射。事实上，具有相同幂零根的理想给出 $\Spec R$ 的同一个拓扑子空间。

但这\emph{是特性而非缺陷}。在例 \ref{ex:0.2.6} 中，理想 $(y,y-x^2) = (x^2,y)$ 不是根理想，它携带的几何含义无法从它的幂零根 $(x,y)$ 看出。具体地说，$f \in (x,y)$ 当且仅当 $f(0,0) = 0$；而 $f \in (x^2,y)$ 当且仅当 $f(0,0) = \partial_x f(0,0) = 0$。粗略地说，这表明 $(y,y-x^2)$ 记住了曲线 $Z(y)$ 与 $Z(y-x^2)$ 在点 $(0,0) \in \A^2_k$ 处彼此相切，且切向量为 $\partial_x|_{(0,0)}$。还应注意，$k[x,y]/(y,y-x^2)$ 的长度等于 $2$，这正是 Bézout 定理所预言的交点个数。

\begin{yizhu}
「长度等于 $2$」是说该环作为 $k$-向量空间的维数为 $2$，等价于说它作为 $\A^2_k$ 中闭子概形的局部环，其极大理想的平方为零而本身非零。原文把它与 Bézout 定理并置，是在为「非既约结构记录交重数」作动员：两条曲线的交点数在经典意义下只能是 $1$（切点），而长度记下了它其实是一阶接触。这一段是全书唯一一处正面论证「为何要允许幂零元」，后文 §\ref{sec:19} 的既约子概形、§\ref{sec:22} 的诺特条件都会回照此处。
[依据：原书 p.7 第 3--4 段；「交重数」的经典表述见 Hartshorne, \emph{Algebraic Geometry}, GTM 52, Ch.~I, \S7（Bézout 定理的陈述）。]
\end{yizhu}

总之，在代数一侧，我们应当考虑所有交换环；在几何一侧，相应的概念称为\emph{仿射概形}。本课程的第一个任务，就是建立如下对偶：
\[
\begin{array}{r|c}
\text{代数} & \text{几何}\\[0.4ex] \hline
\rule{0pt}{1.2em}\text{交换环}\;R & \text{仿射概形}\;X\\
\text{素理想}\;\mathfrak p \subseteq R & \text{点}\;x \in X\\
\text{元素}\;f \in R & \text{函数}\;X \to \A^1_\Z\\
\text{理想}\;I \subseteq R & \text{闭子概形}\;Z \subseteq X.
\end{array}
\]

\begin{yizhu}
末行与第一张表（第 \ref{sec:0.1} 小节）的末行相对照，差别的全部内容就是「根理想 $\to$ 理想」与「Zariski 闭子集 $\to$ 闭子概形」。前者是集合，后者带结构层；例 \ref{ex:0.2.6} 说明省去这层结构就会丢掉切触信息。因此这里的对应必须写成「概形」而不能写成「闭子集」。
[依据：原书 pp.6--7 两张对照表；闭子概形的定义见原书 \S\ref{sec:11.3}。]
\end{yizhu}

% ---------------------------------------------------------------- p.7--8
\subsection{作为结构空间的概形}\label{sec:0.3}

理论上，可以把仿射概形之间的态射定义为由环同态给出的连续映射；再用这类态射把仿射概形粘起来，就得到一般的概形。这模仿了微分流形的定义：概形将是一个配有极大仿射图集的拓扑空间。

实践中，上述做法不便操作。我们更愿意用一种更有效的方式，把附加结构编码到概形底层的拓扑空间上。实现这一点的一个极简单却有力的办法，是拓扑空间上的\emph{层}这一概念；它或许由 Serre 发现，并由 Grothendieck 推广。

粗略地说，$X$ 上的层 $\mathcal F$ 是一个反变对应 $U \mapsto \mathcal F(U)$，把开子集 $U \subseteq X$ 送到某种结构（如集合、群、环）$\mathcal F(U)$，并满足一定的粘接条件。这里「反变」是指：当 $U \subseteq V$ 时，应给出保持既定结构的映射 $\mathcal F(V) \to \mathcal F(U)$。

\begin{example}\label{ex:0.3.1}
设 $X$ 是拓扑空间。把 $U \subseteq X$ 送到 $U$ 上连续函数之环的对应 $U \mapsto C(U,\R)$，是 $X$ 上的交换环层。
\end{example}

类似地，对光滑流形 $X$，$U \mapsto C^\infty(U,\R)$ 是 $X$ 上的交换环层。这促使我们给出：

\begin{predefinition}\label{predef:0.3.2}
\emph{概形}是指一个拓扑空间 $X$，连同其上的交换环层 $\cO_X$，使得 $X$ 局部同构于仿射概形。
\end{predefinition}

这里对开子集 $U \subseteq X$，$\cO_X(U)$ 应当是 $U$ 上的代数函数之环，但后一概念我们尚未定义。不过对仿射概形 $X \simeq \Spec R$，前面的讨论提示应有 $\cO_X(X) \simeq R$。我们将在后续课程中看到，可以由这一点出发反推，得到整个层 $\cO_X$ 的定义。

\begin{yizhu}
「局部同构于仿射概形」的严格含义是：每一点 $x \in X$ 都有开邻域 $U$，使 $(U, \cO_X|_U)$ 同构于某个 $(\Spec R, \cO_{\Spec R})$。所谓「同构」是局部环化空间的同构，即层同构，且在每个茎上诱导的映射把极大理想拉回极大理想。原文把它称为 Pre-Definition 是诚实的：此处尚未定义仿射概形的结构层，也尚未定义环化空间；定理只能等到第~\ref{sec:4}、\ref{sec:5}~节才补齐。
[来源：Stacks Project, Tag 01IJ（Definition 26.9.1：概形是局部同构于仿射概形的局部环化空间，其态射即局部环化空间的态射），\url{https://stacks.math.columbia.edu/tag/01IJ}（2026-11-17 访问）。]
\end{yizhu}

本课程的目标，是定义概形并研究其基本性质。
